
\subsubsection{2.1 Sliding Window and Sub-Quadratic Attention}

Longformer \citep{Beltagy2020Longformer} establishes sliding window attention as a foundational efficiency primitive for long documents, reducing attention complexity from O($n^{2})$ to O($n\cdot w)$ per layer. Mistral 7B \citep{Jiang2023Mistral} applies window-based attention at training time with w = 4096, demonstrating that sub-quadratic attention is compatible with state-of-the-art performance at 7B scale.

Post-hoc conversion from full attention to SWA without retraining is the problem addressed by SWAA \citep{Yu2025SWAA}, which is the closest prior work to our setting. SWAA confirms that naive zero-shot full-model conversion causes catastrophic quality degradation, and proposes a fine-tuning recovery step. SWAT \citep{Fu2025Sliding} examines SWA at training time with proper position encoding adjustments. SWARR \citep{Liu2026Architecture} shows that even supervised fine-tuning is insufficient post-SWA conversion and proposes reinforcement learning for recovery. These works collectively establish that the challenge of zero-shot SWA conversion is well-documented and not trivially solvable.

\textit{Our positioning:} We share the goal of SWA conversion from SWAA but target the zero-shot, no-fine-tuning regime and apply selective rather than full-model conversion. Our central question --- whether selective entropy-guided conversion of 4 of 32 layers avoids the catastrophic failure mode of full-model conversion --- is empirically answered in this paper.

\subsubsection{2.2 Entropy-Based Attention Analysis}

Entropy has been applied to attention heads to identify redundancy and pruning candidates. Michel et al. [2019] demonstrate that a large fraction of attention heads can be removed at test time without significant accuracy degradation, establishing that transformer models contain substantial head-level redundancy. HIES \citep{Choi2025Entropy} combines head importance scores with entropy for structured attention head pruning, achieving quality improvements by using entropy to regularize importance-based selection. Entropy-Lens \citep{Ali2025Entropy} analyzes per-layer entropy profiles as information-processing signatures across transformer layers, finding systematic entropy evolution with depth.

\textit{Our positioning:} These methods operate at the \textbf{head level} for \textbf{pruning} --- they remove heads while preserving layer structure. We operate at the \textbf{layer level} for \textbf{SWA conversion} --- we replace the attention mechanism's effective receptive field while preserving the layer's residual connection, MLP, and parameter count. The aggregation method for layer-level entropy (head-mean vs. head-max) has not been systematically examined in the head-level pruning literature; our ablation finds a 32\% difference in measured concentration between the two methods.

\subsubsection{2.3 Fixed and Local Attention Patterns}

Evidence that transformer attention is not uniformly global provides the conceptual basis for selective conversion. Raganato et al. [2020] demonstrate that fixed positional patterns (diagonal, BOS-attending, EOS-attending) emerge in encoder models and are functionally separable from content-adaptive heads. StreamingLLM \citep{Xiao2024StreamingLLM} exploits sink-token attention patterns for streaming inference with a local + attention-sink window, demonstrating that LLM attention has structure compatible with local constraints.

\textit{Our positioning:} These works provide conceptual support for the existence of layers where global attention is not essential. We operationalize this insight as a per-layer entropy score with explicit stability validation and directly test whether entropy-selected layers tolerate SWA conversion without quality loss.

\subsubsection{2.4 Calibration-Based Model Analysis}

Calibration sets of 100 sequences are standard in quantization methods (GPTQ \citep{Frantar2023}, AWQ \citep{Lin2024}) for collecting activation statistics. Our use of a 100-sequence WikiText-103 calibration set for entropy scoring follows this precedent, adapting it from quantization to layer characterization for SWA conversion.

